% Generated mechanically from frozen input inputs/p01/ja/conventions.tex.
% Do not edit this generated file; edit the bound locale source instead.

% OK, start here
%

\setcounter{chapter}{1}
\chapter{規約}



\phantomsection
\label{conventions-section-phantom}


\section{備考}
\label{conventions-section-comments}

\noindent
これらの文書を執筆する際の規約についての方針は、実際にうまく機能する規約を
選ぶ、というものである。

\section{集合論}
\label{conventions-section-sets}

\noindent
選択公理を伴う Zermelo--Fraenkel 集合論を用いる。\cite{Kunen} を参照されたい。
宇宙は用いない（この点は SGA4 と異なる）。集合論上の問題を強調することは
しないが、すべてが（もちろん）正しいことは確認しており、したがってそれらを
無視するわけでもない。


\section{圏}
\label{conventions-section-categories}

\noindent
圏 $\mathcal{C}$ は対象の集合と、対象の各対に対して与えられる、その間の射の
集合とからなる。言い換えれば、他の文献で「小さい」圏と呼ばれるものである。
「大きな」圏（対象が真の類をなす圏）も用いるが、それは『圏』注記
\ref{categories-remark-big-categories} に列挙されたものに限る。

\section{代数}
\label{conventions-section-algebra}

\noindent
本書では、環とは単位元 $1$ をもつ可換環をいう。したがって環の圏は始対象
$\mathbf{Z}$ と終対象 $\{0\}$ をもつ（後者は $1 = 0$ となる唯一の環である）。
加群は単位的であると仮定する。\cite{Eisenbud} を参照されたい。

\section{記法}
\label{conventions-section-notation}

\noindent
自然数は $\mathbf{N} = \{1, 2, 3, \ldots\}$ の元とする。
整数は $\mathbf{Z} = \{\ldots, -2, -1, 0, 1, 2, \ldots\}$ の元とする。
有理数体を $\mathbf{Q}$ で表す。
実数体を $\mathbf{R}$ で表す。
複素数体を $\mathbf{C}$ で表す。
